With the metric tensor at hand, we may express the \textsc{Hodge}--\textsc{Laplace} operator with respect to hyperspherical coördinates.
For zero-forms, we know the operator to be $\codee\dee f = -\star\dee\!\star\!\dee f$.
We shall denote for brevity $\dee V = \sqrt{\det{(\metrictensor)}}\afrak$, such that $\star\dee f = \sqrt{\det{(\metrictensor)}} \metricg^{ij} \partial_i f \partial_j \iprod \afrak$.
Clearly $\dee(\partial_j\iprod\afrak) = 0$, meaning
\[
  \dee\!\star\!\dee f = \dee \Big( \sqrt{\det{(\metrictensor)}} \metricg^{ij} \partial_i f \Big) \wedge \partial_j \iprod \afrak = \partial_k \Big( \sqrt{\det{(\metrictensor)}} \metricg^{ij} \partial_i f \Big) \dee x^k \wedge \partial_j \iprod \afrak.
\]

\noindent
We know $\dee x^k \wedge \partial_j \iprod \afrak = \delta_{j}^{k} \afrak$, removing all terms of index $k \neq j$.
Applying the \textsc{Hodge} star operator once more to the expression above, we get the following expression for the \textsc{Hodge}--\textsc{Laplace} operator in local coördinates.%
\footnote{%
  \cite{nishikawa2002variational} \textsc{Nishikawa}, p.116
}

\begin{equation}\label{eq:Hodge-Laplace-scalar}
  \Delta f = \frac{1}{\sqrt{\det(\metrictensor)}} \frac{\partial}{\partial x^j} \left( \sqrt{\det(\metrictensor)} \metricg^{ij} \frac{\partial f}{\partial x^i} \right)
\end{equation}

\noindent
It is known that the \textsc{Hodge}--\textsc{Laplace} operator on $\euclidean^{d+1}$ may be decomposed so as to separate the dependence of the radial distance from the origin, and the placement of the point of evaluation on the unit hypersphere \gls{dsphere}.%
\footnote{%
  \cite{chavel1984eigenvalues} \textsc{Chavel}, pp.33--36
}
%
This is a digestible result once one considers the projection of a point from the origin, through the unit hypersphere, unto the designated radius.
This is obviously something we can show using equation \eqref{eq:hyperspherical-metric-tensor}, but shall not do.
Thus we have the following representation the \textsc{Hodge}--\textsc{Laplace} operator.

\begin{equation}\label{eq:hyperspherical-laplacian}
  \Delta_{\euclidean^{d+1}} f = r^{-d} \partial_r \big( r^d \partial_r f \big) + r^{-2} \Delta_{\dsphere^d} f \big( \dsphere(r) \big);
\end{equation}

\begin{equation}\label{eq:angular-hyperspherical-laplacian}
  \Delta_{\dsphere^j} f = \sin^{1-j}{(\alpha^{d-j+1})} \partial_{\alpha^{d-j+1}} \big( \sin^{j-1}{(\alpha^{d-j+1})} \partial_{\alpha^{d-j+1}} f \big) + \csc^{2}{(\alpha^{d-j+1})} \Delta_{\dsphere^{j-1}} f.
\end{equation}

\noindent
Here we have used the recurrence relation of the submatrix determinants of equation \subeqref{\ref{eq:hyperspherical-metric-recurrence}}[b].
To obtain the eigenfunctions of the \textsc{Laplace} equation, we seek solutions to the \textsc{Helmholtz} equation

\begin{subequations}\label{eq:Helmholtz-equation}
  \begin{equation}\tag{\theequation{a, b}}
    \Delta_{\euclidean^{d+1}}\Phi = -\kappa^2 \Phi, \qquad \Phi \in L^2(\euclidean^{d+1}).
  \end{equation}
\end{subequations}

\noindent
From equation \eqref{eq:hyperspherical-laplacian}, we find through multiplying by $r^2$ the confinement of $\kappa$ to the radial part of the equation, perhaps lending some credence to an eigenvalue problem in the angular part separately.
Indeed, since the unit hypersphere is compact, the angular part of the equation has a discrete spectrum.
According to the spectral theorem,%
\footnote{%
  \cite{grigoryan2009heat} \textsc{Grigor'jan}, thm.10.13; thm.10.20
}
the resolvent operator ${(\mathrm{id} - \Delta_{\dsphere^d})}^{-1}$ is compact and selfadjoint,%
\footnote{%
  \cite{conway2007course} \textsc{Conway}, prop.3.9; prop.4.13; thm.5.1
}
with a countable spectral decomposition and the same eigenspace as the angular part of the hyperspherical \textsc{Hodge}--\textsc{Laplace} operator, yielding orthonormal basis with which to expand the operator in.
That is, the full \textsc{Hodge}--\textsc{Laplace} operator preserves the eigenspace of the angular part.
We may then construct the eigenspace iteratively by considering ever higherdimensional unit hyperspheres on which we solve the eigenvalue problem.
For $\alpha^0$, the operator is simply the second derivative, and the eigenvalue problem is the simple harmonic oscillator,

\begin{equation}
  \partial_{\alpha^0}^2 \spherical_{m_0} = - {(m_{0})}^2 \spherical_{m_0}, \qquad \spherical_{m_0}(\alpha^0) = \exp{(i m_{0} \alpha^0)}.
\end{equation}

\noindent
By the construction of the hyperspherical coördinates, in the sense that $\dsphere^1 \subset \dots \subset \dsphere^d$, the ever hyperior hyperspherical harmonics, so to speak, are eigenfunctions for the lowerdimensional spaces as well.
For higherdimensional operators, we write $\spherical_{m_{j-1}}^{m_{j-2}, \dots, m_0}$.
By equation \eqref{eq:angular-hyperspherical-laplacian}, we get the following \textsc{Sturm}--\textsc{Liouville} problem for $j = 2$ when substituting $x = \cos{\alpha^{d-1}}$.

\begin{equation}
  (1 - x^2) F^{\prime\prime} - 2x F^{\prime} + \left( \lambda - \frac{{(m_0)}^2}{1 - x^2} \right) F = 0,
\end{equation}

\noindent
where $\spherical_{m_1}^{m_0} = F \spherical_{m_0}$.
This is the associated \textsc{Legendre} equation, obtained by a substitution in the hypergeometric equation \eqref{eq:hypergeometric}.
We force regular solutions at $x = \pm 1$, yielding the eigenvalues $\lambda = m_1(m_1 + 1)$, such that $m_1 = \absl{m_0} + k$ for all nonegative integers $k$.
In general, we have that the angular hyperspherical eigenvalue problem through equation \eqref{eq:angular-hyperspherical-laplacian} takes the form of a \textsc{Gegenbauer} equation, with solutions in the \textsc{Gegenbauer} polynomials times a sinusoidal function.%
\footnote{%
  \cite{avery2018hyperspherical} \textsc{Avery} \& \textsc{Avery}, \S3.5
}
It is customary to favor simple idices over the index chain for the hyperspherical harmonics, so that the upper index is $m$, and the lower index is $n$.
We seek solutions in threedimensional Euclidean space, meaning we want the spherical harmonics on the two-sphere, which are%
\footnote{%
  \cite{cheng1999fast} \textsc{Cheng}, \textsc{Greengard}, \& \textsc{Rohlin}  (1999), eq.4
}\textsuperscript{,}%
\footnote{%
  \cite{gumerov2003recursions} \textsc{Gumerov} \& \textsc{Duraiswami} (2003), eq.2.5
}

\begin{subequations}\label{eq:sphericalharmonics}
  \begin{equation}\tag{\theequation{a, b}}
    \spherical_{n}^{m}(\alpha^{0}, \alpha^{1}) = c \legendre_{n}^{\absl{m}}{(\cos{\alpha^1})} \exp{(im\alpha^0)}, \qquad c = {(-1)}^m \sqrt{\frac{2n+1}{4\pi} \frac{(n - \absl{m})!}{(n + \absl{m})!}}
  \end{equation}
\end{subequations}

\noindent
These constitute a complete orthonormal basis when $n \geq \absl{m}$.%
\footnote{%
  \cite{takeuchi1994modern} \textsc{Takeuchi}, p.223
}
